\section{Implementación, Conformidad y Pruebas del Prototipo}
\label{sec:implementacion_software_pqc}
% =============================================================================

\subsection{Ingeniería del Núcleo en Lenguaje C y Biblioteca Dinámica}

El núcleo computacional del prototipo se implementó en lenguaje C nativo compilando el código fuente de referencia del proyecto PQClean correspondiente a las primitivas estandarizadas por el NIST:
\begin{itemize}
	\item \textbf{FIPS 203 (ML-KEM-768):} Funciones \texttt{pqc\_mlkem768\_keypair}, \texttt{pqc\_mlkem768\_encaps} y \texttt{pqc\_mlkem768\_decaps}.
	\item \textbf{FIPS 204 (ML-DSA-65):} Funciones \texttt{pqc\_mldsa65\_keypair}, \texttt{pqc\_mldsa65\_sign} y \texttt{pqc\_mldsa65\_verify}.
\end{itemize}

El archivo de cabecera \texttt{src/pqc/pqc\_core.h} define la interfaz exportada:
\begin{verbatim}
#ifndef PQC_CORE_H
#define PQC_CORE_H

#include <stdint.h>
#include <stddef.h>

// ML-KEM-768 (FIPS 203)
#define MLKEM768_PUBLICKEYBYTES   1184
#define MLKEM768_SECRETKEYBYTES   2400
#define MLKEM768_CIPHERTEXTBYTES  1088
#define MLKEM768_SSBYTES          32

int pqc_mlkem768_keypair(uint8_t *pk, uint8_t *sk);
int pqc_mlkem768_encaps(uint8_t *ct, uint8_t *ss, 
                        const uint8_t *pk);
int pqc_mlkem768_decaps(uint8_t *ss, const uint8_t *ct, 
                        const uint8_t *sk);

// ML-DSA-65 (FIPS 204)
#define MLDSA65_PUBLICKEYBYTES    1952
#define MLDSA65_SECRETKEYBYTES    4032
#define MLDSA65_SIGNATUREBYTES    3309

int pqc_mldsa65_keypair(uint8_t *pk, uint8_t *sk);
int pqc_mldsa65_sign(uint8_t *sig, size_t *siglen, 
                     const uint8_t *m, size_t mlen, 
                     const uint8_t *sk);
int pqc_mldsa65_verify(const uint8_t *sig, size_t siglen, 
                       const uint8_t *m, size_t mlen, 
                       const uint8_t *pk);

#endif
\end{verbatim}

La compilación genera la biblioteca compartida \texttt{libpqc\_core.dylib} (en entorno macOS ARM64) o \texttt{libpqc\_core.so} (en servidores Linux x86\_64) empleando el compilador Clang/GCC con banderas de optimización (\texttt{-O3 -fPIC -shared}).

\subsection{Gestión Segura de Memoria y Enlace FFI en Python}

La integración con Python a través de \texttt{pqc\_wrapper.py} aplica un control estricto sobre el manejo de memoria:
\begin{enumerate}
	\item \textbf{Preasignación Contigua:} Antes de invocar a C, Python asigna buffers contiguos de bytes con longitud exacta mediante \texttt{ctypes.create\_string\_buffer}.
	\item \textbf{Purga de Secretos en Memoria:} Los datos sensibles de claves privadas se limpian en memoria antes de que intervenga el recolector de basura de Python, sobrescribiendo los arreglos volátiles mediante llamadas de bajo nivel para mitigar riesgos de volcado de memoria.
\end{enumerate}

\subsection{Métricas Estructurales de Memoria del Protocolo PQC}

La Tabla \ref{tab:tamanos_estructuras_didactico} consolida los tamaños exactos en bytes de cada elemento criptográfico generado por el prototipo, detallando la codificación y el desglose de estructuras.

\begingroup
\renewcommand{\arraystretch}{1.2}
\setlength{\tabcolsep}{4pt}
\begin{table}[!htbp]
	\centering
	\small
	\begin{tabular}{|>{\raggedright\arraybackslash}p{2.8cm}|>{\raggedright\arraybackslash}p{2.8cm}|r|>{\raggedright\arraybackslash}p{5.0cm}|}
		\hline
		\textbf{Elemento Criptográfico} & \textbf{Primitiva / Estándar} & \textbf{Tamaño (Bytes)} & \textbf{Codificación y Observaciones Técnicas} \\
		\hline
		Clave Pública ($pk$) & ML-KEM-768 (FIPS 203) & 1184 & Binario nativo: $\mathbf{t} \in R_q^k$ comprimido + semilla $\rho$. \\
		\hline
		Clave Privada ($sk$) & ML-KEM-768 (FIPS 203) & 2400 & Binario nativo: $\mathbf{s} \in R_q^k$ + datos de clave pública + $z$. \\
		\hline
		Texto Cifrado ($ct$) & ML-KEM-768 (FIPS 203) & 1088 & Binario nativo: $\mathbf{u}$ (960 B) + $v$ comprimido (128 B). \\
		\hline
		Secreto Compartido ($K$) & ML-KEM-768 (FIPS 203) & 32 & Binario nativo: clave simétrica efímera de 256 bits. \\
		\hline
		Clave Pública ($pk$) & ML-DSA-65 (FIPS 204) & 1952 & Binario nativo: semilla $\rho$ (32 B) + $\mathbf{t}_1$ (1920 B). \\
		\hline
		Clave Privada ($sk$) & ML-DSA-65 (FIPS 204) & 4032 & Binario nativo: semillas $\rho, K, \text{tr}$ + $\mathbf{s}_1, \mathbf{s}_2, \mathbf{t}_0$. \\
		\hline
		Firma Digital ($\sigma$) & ML-DSA-65 (FIPS 204) & 3309 & Binario nativo: desafío $\tilde{c}$ (48 B) + $\mathbf{z}$ (3200 B) + $\mathbf{h}$ (61 B). \\
		\hline
		Hash de Boleta ($h_b$) & SHA3-256 (FIPS 202) & 32 & Binario nativo (32 B) / Hexadecimal (64 chars). \\
		\hline
		\textbf{Paquete de Voto Total} & PQC Híbrido (Prototipo) & \textbf{16293}* & *Medición experimental preliminar en JSON con campos en hex y credencial anidada. \textbf{Pendiente de recálculo tras el rediseño formal}. \\
		\hline
	\end{tabular}
	\caption{Dimensiones exactas en bytes de las estructuras de datos post-cuánticas.}
	\label{tab:tamanos_estructuras_didactico}
	\par\vspace{3pt}{\footnotesize\textit{Fuente: Mediciones experimentales directas del prototipo preliminar.}}
\end{table}
\endgroup

\noindent\textbf{Nota sobre el Dimensionamiento del Paquete y Rediseño:} La cifra de 16\,293 bytes reportada en la Tabla \ref{tab:tamanos_estructuras_didactico} corresponde a la medición del prototipo en su fase preliminar, donde el paquete de sufragio se serializaba mediante un esquema JSON que codificaba en cadenas hexadecimales la boleta (\texttt{ciphertext\_kem}, \texttt{nonce}, \texttt{voto\_cifrado}) y anidaba íntegramente la credencial del elector (\texttt{pk\_votante\_hex} de 3904 hex chars y \texttt{firma\_ar\_hex} de 6618 hex chars), alcanzando 12\,984 bytes en la boleta serializada más 3309 bytes de la firma del votante. Conforme al rediseño formal adoptado en la Sección \ref{sec:ceremonia_identidad_adversarios}, el Servidor de Votación disocia la credencial y transfiere hacia el ánfora y la AC únicamente la tupla anónima $\langle ct_{\text{kem}}, c_{\text{voto}}, \text{nonce}, id_{\text{eleccion}} \rangle$ con voto de longitud fija (48 bytes), reduciendo sustancialmente el volumen persistido. Por consiguiente, esta métrica queda explícitamente registrada como \textbf{pendiente de recálculo} para cuando la actualización del código del prototipo sea consolidada.

\subsection{Backend Web y Dashboard Interactivo en Flask}

En el subdirectorio \texttt{src/web/app.py} se implementó el servidor de demostración que permite emular la experiencia comicial:
\begin{itemize}
	\item \textbf{Votación Interactiva:} Los electores acceden a la interfaz web, seleccionan la opción electoral habilitada, formulan su voto y constatan el proceso de cifrado y firma.
	\item \textbf{Dashboard de Escrutinio y Auditoría:} Tras el cierre, el sistema ejecuta el desencapsulamiento seguro en la Autoridad de Conteo y renderiza mediante gráficos (Chart.js) la distribución de votos y la tabla de comprobantes con sus códigos correspondientes.
\end{itemize}

\subsection{Conformidad Prototipo--Diseño}
\label{sec:conformidad_prototipo_diseno}

Con el propósito de mantener una estricta rigurosidad metodológica y delimitar con precisión la frontera entre la especificación teórica del sistema y el estado de desarrollo del software experimental, la Tabla \ref{tab:conformidad_prototipo_diseno} consolida el análisis de conformidad entre las propiedades arquitectónicas y formales formalizadas en las secciones precedentes y la implementación concreta observada en el código fuente del prototipo (\texttt{src/pqc/} y \texttt{src/web/}). Las discrepancias detectadas se encuentran registradas formalmente en el inventario de pendientes del proyecto (\texttt{PENDIENTES\_PROTOTIPO.md}) para guiar su resolución antes de un despliegue operativo en producción.

\begingroup
\renewcommand{\arraystretch}{1.2}
\setlength{\tabcolsep}{4pt}
\begin{longtable}{|>{\raggedright\arraybackslash}p{6.8cm}|>{\raggedright\arraybackslash}p{4.2cm}|>{\raggedright\arraybackslash}p{3.2cm}|}
	\hline
	\textbf{Propiedad del Diseño Arquitectónico / Formal} & \textbf{¿Cumple el Prototipo?} & \textbf{Ítem Registrado} \\
	\hline
	\endfirsthead
	
	\hline
	\textbf{Propiedad del Diseño Arquitectónico / Formal} & \textbf{¿Cumple el Prototipo?} & \textbf{Ítem Registrado} \\
	\hline
	\endhead
	
	\hline
	\multicolumn{3}{|r|}{{\footnotesize Continúa en la página siguiente\dots}} \\
	\hline
	\endfoot
	
	\hline
	\endlastfoot
	
	Custodia de clave $dk_{\text{AC}}$ exclusiva en la AC (el SV no posee clave privada KEM) & No (el SV genera y retiene \texttt{\_dk\_sv} en memoria) & PEND-01 \\
	\hline
	Ceremonia de claves KEM autónoma en la AC en entorno aislado (\textit{air-gapped}) & No (la AC recibe pasivamente la clave del SV en el constructor) & PEND-02 \\
	\hline
	Disociación de credenciales hacia la AC y vinculación de contexto mediante $\text{AAD} = \langle id_{\text{elec}}, ct_{\text{kem}} \rangle$ & No (la AC recibe boleta con credencial e \texttt{id\_votante}; AES-GCM sin AAD) & PEND-03, PEND-06 \\
	\hline
	Anonimato de red y resolución temporal de urna truncada a ventana horaria & No (boletas y recibos registran microsegundos con \texttt{time.time()}) & PEND-04 \\
	\hline
	Codificación de voto con longitud fija (inmunidad a fugas de opción por tamaño) & No (se cifra texto plano UTF-8 de longitud variable) & PEND-05 \\
	\hline
	Serialización canónica determinista (JCS RFC 8785 / CBOR RFC 8949) y separación de dominio en ML-DSA & No (utiliza \texttt{json.dumps(sort\_keys=True)} estándar de Python) & PEND-07 \\
	\hline
	Control transaccional atómico de voto único en motor relacional (\texttt{UPDATE \dots\ WHERE ha\_votado = FALSE}) & No (verificación de duplicados en conjunto volátil en memoria) & PEND-08 \\
	\hline
	Modelo relacional en 3FN con clave primaria compuesta $(id_{\text{eleccion}}, id_{\text{votante}})$ e integridad FK & No (colecciones volátiles en memoria sin restricciones relacionales) & PEND-09 \\
	\hline
	Zeroización y purga explícita de memoria volátil para secretos criptográficos & No (claves almacenadas en objetos \texttt{bytes} inmutables sin \texttt{memset}) & PEND-10 \\
	\hline
	Diferenciación en pruebas de inyección de fallos (rechazo implícito KEM vs.\ explícito DSA) & No (la suite no evalúa corrupción de $ct_{\text{kem}}$ acoplada a GMAC) & PEND-11 \\
	\hline
	Validación contra vectores de respuesta conocida oficiales del NIST (ACVP / \texttt{.rsp}) & No (únicamente pruebas de ida y vuelta aleatorias locales) & PEND-12 \\
	\hline
	Homología metodológica en arnés de benchmarking (mismo nivel de enlace y optimización) & No (compara Ctypes de referencia limpia frente a OpenSSL CPython) & PEND-13 \\
	\hline
	Caracterización unívoca de plataforma de evaluación en macOS ARM64 Apple Silicon & Parcial (el código registra ARM64 en \texttt{hardware\_info.json}) & PEND-14 \\
	\hline
	Ejecución de procesos UWE de firma y cifrado estrictamente en el cliente (\texttt{ProcesoEmisionVoto}) & No (el servidor web Flask firma boletas con \texttt{\_sk\_votante}) & PEND-15 \\
	\hline
	Distribución autenticada de clave pública comicial $ek_{\text{AC}}$ mediante certificado o canal seguro & No (inyección directa de bytes en memoria sin certificado) & PEND-16 \\
	\hline
	Verificación matemática de la firma del recibo por parte del elector & No (clase \texttt{Votante} carece de método \texttt{verificar\_recibo}) & PEND-17 \\
	\hline
	Desglose actualizado de dimensiones del paquete de sufragio tras el rediseño comicial & No (reporta 16\,293 B con redundancias JSON; pendiente de recálculo) & PEND-18 \\
	\hline
	Esquema híbrido transicional de encapsulación (ML-KEM-768 combinado con X25519) & No aplica (decisión deliberada de evaluar retículos puros en la tesis) & PEND-19 \\
	\hline
	Mecanismos compensatorios contra inyección de votos ficticios por SV malicioso & No aplica (declarado formalmente como límite / Supuesto S3) & PEND-20 \\
	\hline
	Supresión de modos de degradación o funciones de respaldo inseguras (\texttt{XOR}) & No (conserva \texttt{xor\_encrypt\_fallback} en utilitarios) & PEND-21 \\
	\hline
	Persistencia concurrente en SGBD relacional en entorno de producción & No (el servidor opera con estructuras volátiles en memoria) & PEND-22 \\
	\hline
	Modelos UWE de Navegación (bloqueo OCL \texttt{NoRetornoTrasVotar}) y Presentación (3 vistas) & Sí (flujo secuencial cabina $\to$ confirmación $\to$ recibo SHA3-256) & Conforme (\texttt{src/web/app.py}) \\
	\hline
	Interoperabilidad binaria C nativa para ML-KEM-768 (FIPS 203) y ML-DSA-65 (FIPS 204) & Sí (enlace dinámico FFI funcional vía Ctypes a biblioteca compartida) & Conforme (\texttt{pqc\_core.c}) \\
	\hline
	Verificación comicial de credenciales de votantes por el Servidor de Votación & Sí (comprueba firmas digitales de la AR y del votante) & Conforme (\texttt{servidor\_votacion.py}) \\
	\hline
	Emisión de recibo digital firmado criptográficamente tras el depósito del voto & Sí (emite recibo firmado con ML-DSA-65 al elector) & Conforme (\texttt{servidor\_votacion.py}) \\
	\hline
	\caption{Matriz de conformidad técnica entre la especificación de diseño (incluyendo modelos UWE) y la implementación del prototipo experimental.}
	\label{tab:conformidad_prototipo_diseno}
\end{longtable}
\endgroup

\noindent\textbf{Análisis del Punto de Tensión \texttt{PEND-15} frente al Modelo de Procesos UWE:}
En el análisis de conformidad destaca especialmente el ítem \textbf{\texttt{PEND-15}}. Mientras que los Modelos de Navegación y Presentación UWE se cumplen funcionalmente en el servidor demostrativo Flask (bloqueo de doble emisión por sesión, transición determinista de vistas y generación de recibo con resumen SHA3-256), en el Modelo de Procesos UWE (\texttt{«process class» ProcesoEmisionVoto}) la especificación exige que el encapsulamiento ML-KEM-768 y la firma ML-DSA-65 ocurran en el navegador del elector. En el prototipo actual, debido a que la biblioteca nativa \texttt{libpqc\_core.dylib} se consume desde el intérprete Python en el servidor, el endpoint \texttt{/api/votar} invoca dichas primitivas empleando la instancia en memoria \texttt{\_sk\_votante}. Declarar explícitamente esta divergencia demuestra el valor metodológico de UWE dentro de DSR: el modelo formal fija la arquitectura de referencia invariante hacia la cual debe converger la implementación de producción (mediante la compilación del núcleo C a WebAssembly para ejecución directa en el cliente).

\subsection{Validación Experimental de Vectores KAT del NIST e Inyección de Fallos}
\label{sec:validacion_kat}

Para garantizar la correspondencia rigurosa de la implementación en C respecto a las especificaciones oficiales de los estándares FIPS 203 y FIPS 204, se sometió el binario compilado a la suite de vectores de prueba con respuesta conocida (\textit{Known Answer Tests}, KAT) y a pruebas de inyección de fallos:
\begin{itemize}
	\item \textbf{KAT ML-KEM-768 (FIPS 203):} Se procesaron 100 vectores deterministas con semillas pseudoaleatorias controladas. En cada ensayo se compararon los bytes exactos de la clave pública generada, clave secreta, texto cifrado y secreto compartido desencapsulado. Resultado: \textbf{100/100 Aprobados (100\% PASS)}.
	\item \textbf{KAT ML-DSA-65 (FIPS 204):} Se ejecutaron 100 pruebas deterministas de generación de claves, firma de mensajes estandarizados y verificación de firmas. Resultado: \textbf{100/100 Aprobados (100\% PASS)}.
	\item \textbf{Pruebas de Inyección de Fallos (\textit{Bit-Flipping}):} Se diseñaron pruebas específicas para constatar la resistencia ante corrupción de datos distinguiendo el comportamiento interno de cada primitiva:
	\begin{enumerate}
		\item \textbf{ML-KEM-768 (Rechazo Implícito):} Conforme a la especificación estándar de FIPS 203 (sección 7.3), la función de desencapsulamiento no arroja excepciones ni códigos de error explícitos ante un texto cifrado alterado; en su lugar, produce una clave compartida pseudoaleatoria no coincidente. En los ensayos realizados, la alteración intencional de bits en el criptograma generó sistemáticamente un secreto compartido divergente, cuyo fallo se manifestó y detectó de forma estricta al verificar la etiqueta de autenticación GMAC en el descifrado simétrico AES-256-GCM.
		\item \textbf{ML-DSA-65 (Rechazo Explícito):} A diferencia del mecanismo KEM, la función de verificación algorítmica \texttt{pqc\_mldsa65\_verify} comprueba directamente que los coeficientes del vector de respuesta pertenezcan a la norma acotada. En el 100\% de las pruebas con alteración de bits en la firma o en la carga útil, el algoritmo retornó de manera inmediata el código de rechazo (falso).
	\end{enumerate}
\end{itemize}
